\chapter{指南使用说明与研发测试路线图}
\section{本书定位及维护规则}
本文件是\textbf{动态倾角仪全研发周期可持续扩展的测试指南}，不是仅针对某一次常温标定的验收报告，也不是各阶段测试已全部完成的声明。每阶段均遵守\textbf{阶段目标和前置条件 $\rightarrow$ 测试项目总表 $\rightarrow$ 逐项实施方案 $\rightarrow$ 数据记录与阶段评审}的顺序。
\begin{itemize}
\item \textbf{章节导航：} S0、S1、S2、S3、S4 各自独立成章。阶段顺序代表推荐研发工作流，不代表必须一次做完才可以启动下一项。
\item \textbf{编号稳定：} 当前已使用的 CT-00--CT-18 全部保留；新增 Allan 方差标为 \textbf{CT-12A}，不影响此前试验台账与脚本关联。后续新增测试使用本阶段的新编号，不重排历史编号。
\item \textbf{成熟度明确：} S1 项目已有具体测试方案；S0 列出已有工艺的应检项目及方法概要；S2--S4 先建立完整候选测试清单和能力边界，标注\textbf{待细化／待设备验证}。
\item \textbf{结果单独归档：} 同一编号在每次试验中应使用独立 \texttt{run\_id}，记录“未执行、执行、阻断、通过、失败、诊断性结论”，不可把条目列入指南就当作已通过。
\end{itemize}
\section{研发阶段导航总表}
\small
\begin{longtable}{C{13mm}p{30mm}p{64mm}p{36mm}}
\toprule
阶段 & 测试主题 & 本阶段核心任务 & 当前文档状态\\\midrule\endfirsthead
\toprule 阶段 & 测试主题 & 本阶段核心任务 & 当前文档状态\\\midrule\endhead
S0 & 常温内参标定 & ACC/GYRO 原始数据检查、零偏与全矩阵拟合、独立选模、参数冻结与追溯 & 工艺/测试目录 + 要点；需现场复核\\
S1 & \textbf{常温标定后验证} & 标定后残余、静态倾角、复合姿态、陀螺仪、重复性、噪声及 Allan 方差、低速动态诊断 & \textbf{总表 + 20 项详细方案（当前主章节）}\\
S2 & 全温和温度补偿 & 多温度点、升降温迟滞、温度动态、温补参数与全温残余精度 & 测试总表；具体步骤待细化\\
S3 & 动态倾角与抗干扰 & 动态精度/响应、相位和时间同步、线加速度、离心、振动及冲击 & 测试总表；参考系统能力待核实\\
S4 & 可靠性、复测与批量 & 长期漂移、环境试验前后、工装 GR\&R、批次一致性、制造放行 & 测试总表；样本量/标准待确定\\\bottomrule
\end{longtable}
\normalsize
\section{应如何定位 Allan 方差（Allan variance）}
\textbf{已经纳入本指南。} 它不是用来求解确定性零偏和比例因子矩阵的“另一种六面体标定”，而是用于表征\textbf{静态 IMU 随机噪声与零偏稳定性}的分析方法，因此安排在 \textbf{S1 常温标定后验证}。
\begin{itemize}
\item CT-12：短期静态噪声、均值波动和短期漂移，\textbf{P0}，可使用 30--60 分钟记录。
\item \textbf{CT-12A：Allan 方差／Allan 偏差分析}，\textbf{P1}，单独采集更长的连续静态数据并评估角随机游走、零偏不稳定性等随机项。
\item S2、S4 可在不同温度或老化后\emph{重复执行 CT-12A 的方法}，并使用阶段自身的温度/稳定性编号形成对比，但不混淆“常温基线”和“全温/寿命结论”。
\end{itemize}
\section{内参与整机性能的边界}
\begin{itemize}
\item \textbf{传感器层：} 验证 ACC/GYRO 输出的残余偏置、比例、三轴耦合、噪声以及随机稳定性。
\item \textbf{静态倾角层：} 可用补偿后 ACC 和独立静态姿态计算 Roll/Pitch；若 PCB 不输出融合角，不得把离线角度称作动态融合精度。
\item \textbf{机械安装层：} PCB 坐标系 $S$、转台平台参考系 $P$ 和未来整机坐标系 $B$ 不可混同。当前裸 PCB 的装夹外参需要记录。
\end{itemize}
\begin{warningbox}
本阶段设备是一块 IMU PCB 与 2TD-600 双轴转台。说明书未给出\emph{角速率精度保证}，也未声明具备纯线性加速度或振动试验能力。S2/S3 中需要温箱、振动台、独立同步或溯源参考的条目必须写为\textbf{有条件／待配置}，不能将“规划列项”等同“目前能做”。
\end{warningbox}
